
\paragraph{Problem 1(b) answer}\GradescopeComment{This should be on page 3 of your PDF.}

\begin{enumerate}[label=(\roman*)]
\item[] Here is the recurrence relation for $M$.
  \(M(1) = \BLUE{
    %
    \REPLACEME
    % 
  }\).
  %
  For $n\ge 2$, \BLUE{
    \[ M(n) =
      %
      \REPLACEME
      %
    \]
  }

\item[] In level $i$ of the recursion tree, the number of nodes is \BLUE{
    \[
      \REPLACEME
    \]
    }
  
\item[] 

  The size of each subproblem in level $i$ is \BLUE{
    \[
      \REPLACEME 
    \]
  }

\item[] 

  The work associated with each node in level $i$ is \BLUE{
    \[
      \REPLACEME 
    \]
  }

\item[] 

  The total work in level $i$ is \BLUE{
    \[
      \REPLACEME 
    \]
  }

\item[]

  The total number of levels is \BLUE{
    \[
      \REPLACEME 
    \]
  }
  
  By summing over the levels, $M(n)$ is proportional to the following sum:
  
\item[]
  \BLUE{
    \[
      \REPLACEME   %% \sum_{i=\ldots}^{\ldots}\Theta(\ldots)
    \]
  }

\vfill 

\item[]
  
  The big-$\Theta$-value of $M(n)$ is \BLUE{
    \[
      \REPLACEME
    \]
  }
  %  NO EXPLANATION NEEDED
\end{enumerate}


%%% Local Variables:
%%% mode: latex
%%% TeX-master: "homework_2"
%%% End:
